% ===== §2 Background =====
\section{Background and Theoretical Frameworks}
\label{sec:background}

\noindent
The paper draws on, and defines itself against, several established frameworks. This section surveys them and states the paper's position with respect to each, in the order in which they enter the subsequent development. The survey is organised under six headings: joint attention, the grammatical modelling of interaction, the ontology of value, the quantum-structural modelling of cognition, relational and process ontology, and the method of cross-disciplinary unification.

\subsection{Joint attention}
\label{sec:bg-joint}

The concept of joint attention originates in developmental psychology, where it designates the triadic structure in which two subjects orient together toward a shared object, each aware of the other's orientation \citep{tomasello1999,trevarthen1979}. In Tomasello's account the capacity for joint attention is foundational for the acquisition of language and culture, the shared orientation toward a third term being the condition under which symbols acquire their intersubjective stability \citep{tomasello1999}. The phenomenological tradition supplies a distinct but compatible account of intersubjectivity, in which the other is given not as an inferred subject behind a body but in a primordial attunement of embodied subjects: Husserl's analysis of intersubjective constitution \citep{husserl1960}, Merleau-Ponty's account of intercorporeality, in which two bodies are attuned prior to any concept \citep{merleau2012}, and Zahavi's systematic development of the phenomenological position \citep{zahavi1975}.

The present paper retains the concept and the lineage but displaces the structure. The joint attention of the developmental account is triadic and objectual: its third term is an object antecedently present in the world. The co-apprehension of value treated here requires a joint attention whose third term is not an antecedent object but a value in the course of emergence, transversal across the registers and resident in none. This is joint attention at a limit the developmental account does not reach, more adequately described by the phenomenology of intercorporeality and by the psychoanalytic account of an attention organised around a void; the prior paper of this series established the corresponding mechanism at the constitution of the subject, and the present paper displaces it to the constitution of a value between subjects.

\subsection{The grammatical modelling of interaction}
\label{sec:bg-grammar}

The treatment of structured generativity by formal grammars descends from the Chomsky hierarchy, which classifies grammars and their languages by the form of their production rules, from the regular through the context-free and context-sensitive to the recursively enumerable \citep{chomsky1965,hopcroft2006}. The hierarchy and its associated automata supply the standard formal vocabulary for generative structure, and the series has employed probabilistic context-free grammars in this vocabulary to model the recursion of generative justice \citep{eglash2016,eglash1999}. The modelling of interaction specifically, as distinct from the generation of strings, has been pursued in game-theoretic and dialogical formalisms and in the various interaction grammars of computational linguistics, in which the units generated are not sentences but moves in a structured exchange.

The present paper employs the grammatical vocabulary, in the sense that it treats the dynamics of a relation as a generative grammar, the relational grammar of \xref{sec:ontology}. But it argues, in \xref{sec:modelling}, that the relation it treats exhibits a feature no grammar of the Chomsky hierarchy accommodates: the production rules are themselves rewritten, within the process of generation, by what the process generates. This feature is shown to be orthogonal to the hierarchy rather than a position within it, and the consequence for the formal programme is examined at length.

\subsection{The ontology of value}
\label{sec:bg-value}

The paper engages several traditions in the theory of value. In the Marxian tradition, value is not a natural property but a social form, generated in production and circulation, and capable, under the wage relation, of turning against and dominating its producer: the analysis of value, the commodity form, and alienated labour in the \emph{Manuscripts} and in \emph{Capital} \citep{marx1959,marx1976,marx1973}. In the psychoanalytic tradition, value is bound to the economy of desire and enjoyment, organised around the objet a and jouissance of the Lacanian registers \citep{lacan1992,lacan1978sem11,lacan1998,fink1995,evans1996}. In the tradition of generative justice, value must be retained by and returned to those who create it, the criterion of justice lying in the circulation of value through the generative network rather than its extraction from it \citep{eglash2016}. And the phenomenological theory of value, in Scheler's non-formal ethics, treats values as given in feeling, in a stratified order, irreducible to the goods that bear them \citep{scheler1973}.

The present paper does not adopt a realism in which value is an antecedent property awaiting detection, nor a pure constructivism in which value is exhausted by its symbolic positing. It advances, in \xref{sec:registers} and \xref{sec:ontology}, an ontology of value as a phenomenon emergent on the coupled dynamics of the three registers, transversal across them and reducible to none, and it relocates the three registers of value the series has employed, the imaginary, the symbolic, and the real, as aspects internal to this emergent. The Marxian account of the reconfiguration of the producer by the product is taken up, in \xref{sec:political}, as a realisation of the paper's reflexive structure in the register of political economy.

\subsection{The quantum-structural modelling of cognition and the social}
\label{sec:bg-quantum}

The application of the formal structure of quantum theory to cognition and decision, independently of any claim concerning a physical quantum substrate, constitutes the field of quantum cognition. Its results model order effects, interference, and contextuality in judgement and decision by the formal apparatus of state superposition, non-commuting observables, and basis-dependent measurement \citep{busemeyer2012,aerts2009}. The distinction on which the field rests, between the employment of quantum formalism as a mathematical structure and the assertion of a physical quantum substrate, is the distinction the present paper observes. Beyond cognition, the extension of quantum structure to social ontology has been proposed, most expansively by Wendt, who advances a literal quantum social science in which social phenomena are grounded in physical quantum processes \citep{wendt2015}; the present paper does not adopt this literal position and is to be distinguished from it explicitly. The geometric phase, which the series employs to characterise the good cycle, originates in this quantum context, in the phase a state accumulates around a loop in its space of rays \citep{berry1984,pancharatnam1956,aharonov1987}.

The present paper employs quantum structure, in \xref{sec:grammar}, strictly as an epistemic model of participatory generation, with the physical thesis explicitly declined. It diverges from quantum cognition in its object, treating the co-generation of relational value rather than individual judgement, and from literal quantum social science in declining the substrate claim. In \xref{sec:modelling} it considers, and examines the cost of, a more developed quantum-field-theoretic modelling of the dynamics, the assessment of which is one of the paper's principal formal contributions.

\subsection{Relational and process ontology}
\label{sec:bg-process}

The ontology the paper derives is a process ontology and a relational one, and it stands in a lineage. In Whitehead's process philosophy, the fundamental entities are not enduring substances but events or occasions of becoming, and what appears as a thing is an abstraction from a process \citep{whitehead1978}. In Simondon's philosophy of individuation, the individual is not given antecedently but is constituted in a process of individuation that does not exhaust the pre-individual field from which it arises, the individual and its associated milieu being generated together \citep{simondon2020}. In Deleuze, difference and repetition are prior to the identities they generate, and the new is produced rather than reproduced \citep{deleuze1994}. These positions share the priority of process over substance, of generation over the generated, and of relation over the related terms.

The present paper's ontology, in \xref{sec:ontology}, is continuous with this lineage and adds a determinate clause. Relation, grammar, and value are three aspects of one generative system, of which the subject is a product rather than an external premise; and the system's dynamics is reconfigured by its own products, the reflexive clause that distinguishes the paper's process ontology and that conditions the formal dilemma of \xref{sec:modelling}.

\subsection{The method of cross-disciplinary unification}
\label{sec:bg-method}

The paper's final methodological commitment concerns the relation among the disciplinary frameworks it employs. Rather than subordinating them to a single master framework, or treating them as an eclectic assembly, it specifies them as symmetry-broken phases of one abstract structure, related as the realisations of one structure under the conditions of distinct disciplines. The figure of symmetry-breaking is borrowed from physics; the methodological position it supports, that no single framework possesses the whole and that the unbroken structure is cognised only across the phases and not from above, is developed in \xref{sec:breaking} and is there argued to be a dialectical-materialist rather than a relativist commitment. The geometric invariant proposed as conserved across the phases, the holonomy on which the good cycle turns, is advanced as the paper's most speculative conjecture and is left, as the argument requires, unsettled.

With the frameworks surveyed and the paper's position with respect to each stated, the development proceeds, beginning with the three registers as coupled systems.
